Web applications fail in ways that are easy to describe and hard to find: a total that
ignores a changed quantity, a validation rule that holds on one path and not on another, a
value that is lost on reload. Automated exploratory testers can reach such states, but they
face two problems. They must decide where to go in a large, stateful interface where the
decisive actions are many steps deep, and they must prove what they find, or their reports
are noise. This thesis builds and evaluates \emph{Katalyze}, an automated functional
tester for web applications that uses reinforcement learning (RL). It confirms a fault only when the fault
reproduces twice in fresh browser sessions, is absent from a control build without that
fault, and survives delta-debugging minimisation.

The work has two generations. The first (v1) is an action-conditioned dueling double deep
Q-network that scores every available browser action, with prioritised multi-step replay, a
residual head for online adaptation, and a reward shaped by a local large language model
(LLM) that judges task progress. Two pre-registered experiments tested it. Pretraining on
five applications did not speed up fault discovery on a sixth (median difference 0
interactions, 95\% interval $[-749, +290]$). A diagnostic over twelve generated application
variants found that the network does not learn which action helps in which state
(within-label AUROC 0.51). Post-hoc checks showed that the LLM's progress score and its page
labels did not generalise beyond the data they were tuned on.

These findings shaped the second generation (v2), a hybrid tester. It abstracts pages
explicitly, keeps an archive of reached states that it returns to by exact replay
(Go-Explore), and applies test operators from standard test design. A light online learner
(a Thompson-sampling bandit, or count-based curiosity Q-learning) decides which operator to
try. On eight held-out applications, in a pre-registered evaluation of 160 runs with 32
clean-build controls, the hybrid arms verified the most distinct faults (19 and 18 of 24,
against 13 for random exploration) with zero false positives. All four pre-registered
comparisons of time to first fault were inconclusive.
After re-tuning on rebuilt full-stack development applications, the frozen tester was
compared with WebExplor and WebRLED on their six subject applications, measured as they were
measured: one-hour runs, Istanbul branch coverage and console failures. It did not reach
their published numbers. It averaged 44.3\% branch coverage against 55.5\% and 64.2\%, and
1.8 distinct console failures against 3.1 and 4.1. A post-hoc follow-up traced much of the
gap to the tester's action space rather than its policy. With clickable elements that plain
HTML scanning misses, coverage on Phoenix rose from 49\% to 83\% for every arm, level with
WebRLED. The evidence supports
a modest conclusion: the structure of a tester, more than the learning in it, determines
which functional faults it verifies.

\vspace{0.5cm}
\noindent\textbf{Keywords:} automated web testing, reinforcement learning, deep Q-network,
Go-Explore, Thompson sampling, test oracle, seeded faults, survival analysis.
